\begin{table}[h]
	\caption{Indikator Keberhasilan}
\label{tab:indikator keberhasilan}
\begin{tabular}{ | p{1cm} | p{12cm} | }
    \hline
    \textbf{No} & \textbf{Indikator Keberhasilan} \\ \hline

    1 & \textbf{Akurasi Minimal 85\%}: 
    Hasil perhitungan \textit{true positive} dan \textit{true negative} harus menghasilkan skor akurasi di atas 85\%. 
    Angka ini menjadi batas toleransi agar sistem layak digunakan menggantikan pencarian manual. \\ \hline

    2 & \textbf{Waktu Komputasi < 1 Detik/Foto}:
    Sistem harus mampu melakukan deteksi, ekstraksi fitur, dan pencocokan dengan kecepatan rata-rata di bawah 1 detik per foto pada \textit{dataset} pengujian. \\ \hline

    3 & \textbf{Skalabilitas Kapasitas}:
    Sistem mampu memproses \textit{dataset} yang besar foto dalam satu kali eksekusi (\textit{batch processing}) tanpa mengalami kehabisan memori (\textit{Out of Memory}). \\ \hline

    4 & \textbf{Stabilitas Tanpa Crash}:
    Sistem harus memiliki tingkat keberhasilan eksekusi 100\% (tidak \textit{force close}) meskipun terdapat \textit{file} tidak valid atau foto tanpa wajah dalam \textit{dataset} \textit{input}. \\ \hline

\end{tabular}
\end{table}